% Original: PDF strana 2, gornji slajd.
\slajdid{02-s003}
\begin{frame}[label=02-s003]{Kašnjenje invertora pri promeni signala}
% element: 02-s003:e01
\begin{columns}[c,onlytextwidth]
\begin{column}{.72\textwidth}
Signal je fizička veličina koja nosi informaciju o logičkim nulama i jedinicama.
\end{column}
\begin{column}{.24\textwidth}\centering
\begin{tikzpicture}[DEoznaka]
\node[not gate US,draw,minimum width=8mm,minimum height=7mm] (g) at (0,0) {};
\draw[DEvod] (g.input)--++(-.35,0) node[left] {\(x(t)\)};
\draw[DEvod] (g.output)--++(.35,0) node[right] {\(y(t)\)};
\end{tikzpicture}
\end{column}
\end{columns}
\par\vspace{2mm}
% element: 02-s003:e02
\begin{columns}[T,onlytextwidth]
\begin{column}{.43\textwidth}\centering
\Oznaka{Idealizovani odziv}\vspace{3mm}
\begin{tikzpicture}[DEoznaka,x=10mm,y=6.2mm]
\draw[DEstrelica] (-.7,2.55)--(3.8,2.55) node[below] {\(t\)};
\draw[DEstrelica] (0,2.40)--(0,4.35) node[left] {\(x(t)\)};
\draw[DEstrelica] (-.7,-.1)--(3.8,-.1) node[below] {\(t\)};
\draw[DEstrelica] (0,-.25)--(0,1.7) node[left] {\(y(t)\)};
\node[above left,inner sep=1pt] at (0,3.65) {=1};\node[above left,inner sep=1pt] at (0,2.75) {=0};
\node[above left,inner sep=1pt] at (0,1) {=1};\node[above left,inner sep=1pt] at (0,.1) {=0};
\draw[DEvod] (-.5,2.75)--(.6,2.75)--(.6,3.65)--(1.6,3.65)--(1.6,2.75)--(3.5,2.75);
\draw[DEvod] (-.5,1)--(.6,1)--(.6,.1)--(1.6,.1)--(1.6,1)--(3.5,1);
\draw[dashed] (0,3.65)--(.6,3.65);
\draw[dashed] (.6,2.75)--(.6,-.25) node[below] {\(t_1\)};
\draw[dashed] (1.6,2.75)--(1.6,-.25) node[below] {\(t_2\)};
\draw[DEcrvena] (-.8,-.65)--(3,4.15) (-.8,4.15)--(3.5,-.65);
\end{tikzpicture}

\end{column}
\begin{column}{.53\textwidth}\centering
\Oznaka{Odziv sa kašnjenjem}\vspace{3mm}
\begin{tikzpicture}[DEoznaka,x=8.5mm,y=6.2mm]
\draw[DEstrelica] (-.7,2.55)--(3.8,2.55) node[below] {\(t\)};
\draw[DEstrelica] (0,2.40)--(0,4.35) node[left] {\(x(t)\)};
\draw[DEstrelica] (-.7,-.1)--(3.8,-.1) node[below] {\(t\)};
\draw[DEstrelica] (0,-.25)--(0,1.7) node[left] {\(y(t)\)};
\node[above left,inner sep=1pt] at (0,3.65) {=1};\node[above left,inner sep=1pt] at (0,2.75) {=0};
\node[above left,inner sep=1pt] at (0,1) {=1};\node[above left,inner sep=1pt] at (0,.1) {=0};
\draw[DEvod] (-.5,2.75)--(.6,2.75)--(.6,3.65)--(1.6,3.65)--(1.6,2.75)--(3.5,2.75);
\draw[DEvod] (-.5,1)--(1.15,1)--(1.15,.1)--(2.5,.1)--(2.5,1)--(3.5,1);
\draw[dashed] (0,3.65)--(.6,3.65) (0,.1)--(1.15,.1);
\draw[dashed] (.6,2.75)--(.6,-.3) (1.6,2.75)--(1.6,-.3);
\draw[dashed] (1.15,.1)--(1.15,-.3) (2.5,.1)--(2.5,-.3);
\draw[DEcrvena,<->] (.6,-.25)--(1.15,-.25) node[midway,below] {\(t_{pHL}\)};
\draw[DEcrvena,<->] (1.6,-.25)--(2.5,-.25) node[midway,below] {\(t_{pLH}\)};


\draw[DEcrvena,<->] (.6,4.05)--(1.6,4.05) node[midway,above,inner sep=1pt] {\(T_{\mathrm{ul}}\)};
\draw[DEcrvena,<->] (1.15,1.42)--(2.5,1.42) node[midway,above,inner sep=1pt] {\(T_{\mathrm{izl}}\)};
\node[right,inner sep=1pt,text=DEcrvena] at (4.2,2.0) {\(T_{\mathrm{ul}}=T_{\mathrm{izl}}\;???\)};
\end{tikzpicture}

\end{column}
\end{columns}
\par\vspace{2mm}
% element: 02-s003:e03
\begin{columns}[T,onlytextwidth]
\begin{column}{.47\textwidth}
\(\textcolor{DEcrvena}{t_{pHL}}\): kašnjenje pri prelazu izlaza \(H\to L\).
\end{column}
\begin{column}{.49\textwidth}
\(\textcolor{DEcrvena}{t_{pLH}}\): kašnjenje pri prelazu izlaza \(L\to H\).
\end{column}
\end{columns}
\par\vspace{1mm}
\Oznaka{\(p\) — propagation; \(H\) — high, \(H=1\); \(L\) — low, \(L=0\).}
\end{frame}
